% =============================================================================
% conclusion.tex — Conclusion (draft v0.3, 2026-09-03)
% =============================================================================

\section{Conclusion}
\label{sec:conclusion}

We presented a multi-product, market-clearing supply chain model for a
Wisconsin biochar CDR system and embedded three carbon policy architectures in
it. We report three findings. First, for this demand-constrained, net-negative
network, policy design matters more than stringency within the tested ranges:
a linear carbon tax and net
emission caps looser than the baseline flux leave the optimum unchanged---
tightened below it, the cap binds first as an expansion mandate and then as a
500$^\circ$C durability mandate at \$110--330/tCO$_2$e---
while throughput-scaled removal credits activate the bulk
biochar market at \$25/tCO$_2$e under the calibrated step demand and mobilize
essentially all near-term biomass
at \$200/tCO$_2$e, increasing system surplus from \$0.71~B to \$3.36~B/yr and
deepening net removal from 1.60 to 1.96~Mt CO$_2$e/yr. Second, the shadow
marginal abatement cost (fixed-layout LP dual) traced by the cap-and-trade
sweep ($\approx$ \$351--391/t CO$_2$e
over the first binding intervals, rising to $\approx$ \$599/t) lies above
EU ETS allowance prices and the voluntary removal-credit band, quantifying how
expensive gross-emission abatement is for a pyrolysis-based system, and why
removal crediting, not gross-emission pricing, is the effective policy lever.
Third, regulatory detail measurably reshapes
technology: the VM0044 H/C~$\le$~0.7 gate excludes all torrefaction-like
300$^\circ$C biochar from crediting and converts the system to 500$^\circ$C
pyrolysis as credit prices rise (2.25~Mt near-term, 4.25~Mt mature-market at
\$200/tCO$_2$e), and
quality-differentiated demand alone sustains 0.8~Mt/yr of 500$^\circ$C biochar
in the no-policy baseline.

For Wisconsin and the Midwest, the practical message is that feedstock is not
the constraint; demand is. Policies that underwrite farmer adoption or
monetize durable carbon at voluntary-market rates expand the system far more
effectively than supply-side measures, and the spatial structure of value
(biochar inherent values of \$275--324/t across scenarios) implies that
location rents are sensitive to how much biomass the market ultimately calls
for. Extensions under development include seasonal inventory and multi-period
investment, the integration of anaerobic-digestion waste streams that
characterize the dairy economy, and tighter optimality certification on
high-performance computing for the free-facility policy sweeps.
